% Abhakseite „Das kann ich“ – Zone nur im Gesamt (\mitzone)
%% TODO Ich-kann-Sätze: Platzhalter wie in den Aufgabendateien
\begin{abhakseite}
\ifdefined\mitzone
\abhakgruppe{Kennst du schon}
\abhak{1}{Mit Lineal und Geodreieck zeichnen}
\abhak{2}{Addieren und Subtrahieren im Bereich bis zum Vollwinkel, auch mit Dezimalzahlen (Vielfaches minus zwei Dezimalgrade)}
\abhak{3}{Längeneinheiten m und km umrechnen (1,5 km = 1500 m)}
\abhak{4}{Vierecksarten kennen}
\abhak{5}{Kreis mit dem Zirkel zeichnen, Radius und Durchmesser unterscheiden}
\abhak{6}{Addieren und Subtrahieren im Bereich bis zum Vollwinkel, auch mit Dezimalzahlen (Vielfaches minus zwei Dezimalgrade) – Fehler finden}
\abhak{7}{Addieren und Subtrahieren im Bereich bis zum Vollwinkel, auch mit Dezimalzahlen (Vielfaches minus zwei Dezimalgrade) – selbst rechnen}
\fi
\abhakgruppe{Einheit 1 · Winkel messen und zeichnen}
\abhak{8}{Welche Skala?}
\abhak{9}{Winkel messen}
\abhak{10}{Winkel schätzen}
\abhak{11}{Winkel messen – Fehler finden, Begründen, Anwendung}
\abhakgruppe{Einheit 2 · Winkel an Geradenkreuzungen und Parallelen}
\abhak{12}{Sind die Geraden parallel?}
\abhak{13}{Winkel an Geraden}
\abhak{14}{Parallelität aus gleichen Stufenwinkeln prüfen}
\abhak{15}{Winkel an Geraden – Fehler finden, Begründen, Anwendung}
\abhakgruppe{Einheit 3 · Winkelsummen, Dreiecke und Vierecke}
\abhak{16}{Welches Teildreieck?}
\abhak{17}{Winkelsumme}
\abhak{18}{Innenwinkelsumme von Fünf- und Sechseck über Dreiecke}
\abhak{19}{Winkelsumme – Fehler finden, Begründen, Anwendung}
\abhakgruppe{Einheit 4 · Dreiecke konstruieren}
\abhak{20}{Konstruieren}
\abhak{21}{Dreieck skizzieren und beschriften}
\abhak{22}{Konstruieren – Fehler finden, Begründen, Anwendung}
\abhakgruppe{Einheit 5 · Besondere Linien im Dreieck und Satz des Thales}
\abhak{23}{Besondere Linien}
\abhak{24}{Besondere Linien – Fehler finden, Begründen, Anwendung}
\end{abhakseite}
